\section{Utilisation de l'app}

\subsection{Compilation et utiliseation sur la Raspberry Pi}
    Nous utilison la cross compilation pour compiler le code pour la Raspberry Pi.
    
    Une fois les dépendance GCC et Rsync installé, il faut setupe le ".env" à partir du ".env-template" fourni dans le projet.
    Pui il faut utiliser les commandes suivantes:
    \begin{lstlisting}[language=bash]
        make sync-sysroot
        make 
        make remote-copy
    \end{lstlisting}

    Ensuite il faut acceder en SSH à la Raspberry Pi et executer le binaire:
    \begin{lstlisting}[language=bash]
        make remote-ssh
        cd /home/pi/cross-compile
        ./<Nom_de_l_executable_setup_dans_env>
    \end{lstlisting}

    Après l'execution vous pourrez constater les Log dans le dossier log situé dans le dossier d'éxécution.